% ============================================================
%  PREAMBLE — Relazione Progetto Introduzione alla Prog. Web
% ============================================================

% --- Classe documento ---
\documentclass[12pt, a4paper, oneside]{article}

% --- Encoding & lingua ---
\usepackage[T1]{fontenc}
\usepackage[utf8]{inputenc}
\usepackage[italian]{babel}

% --- Font ---
\usepackage{lmodern}
\usepackage[patch=none]{microtype}

% --- Geometria pagina ---
\usepackage[
  top=2.5cm, bottom=2.5cm,
  left=3cm,  right=3cm,
  headheight=15pt
]{geometry}

% --- Intestazione e piè di pagina ---
\usepackage{fancyhdr}
\pagestyle{fancy}
\fancyhf{}
\fancyhead[L]{\small\textit{Relazione Progetto — Introduzione alla Programmazione per il Web}}
\fancyhead[R]{\small\textit{A.A.\ 2025--2026}}
\fancyfoot[C]{\thepage}
\renewcommand{\headrulewidth}{0.4pt}

% --- Colori ---
\usepackage[dvipsnames, table]{xcolor}
\definecolor{primary}{RGB}{30, 80, 160}      % blu accademico
\definecolor{secondary}{RGB}{90, 90, 90}     % grigio testo secondario
\definecolor{codebg}{RGB}{245, 245, 245}     % sfondo listati

% --- Hyperref (ultima, salvo eccezioni) ---
\usepackage[
  colorlinks = true,
  linkcolor  = primary,
  urlcolor   = primary,
  citecolor  = primary,
  pdftitle   = {Relazione Progetto Web},
  pdfauthor  = {Team},
  pdfsubject = {Introduzione alla Programmazione per il Web},
]{hyperref}

% --- Grafici e diagrammi ---
\usepackage{graphicx}
\usepackage{float}
\usepackage{tikz}
\usetikzlibrary{shapes.geometric, arrows.meta, positioning, fit, backgrounds, calc}
\usepackage{pgf-umlsd}          % diagrammi di sequenza (opzionale)

% --- Tabelle ---
\usepackage{booktabs}
\usepackage{tabularx}
\usepackage{multirow}
\usepackage{array}

% --- Listati di codice ---
\usepackage{listings}
\lstset{
  backgroundcolor = \color{codebg},
  basicstyle      = \ttfamily\small,
  keywordstyle    = \color{primary}\bfseries,
  commentstyle    = \color{secondary}\itshape,
  stringstyle     = \color{BrickRed},
  numbers         = left,
  numberstyle     = \tiny\color{secondary},
  stepnumber      = 1,
  numbersep       = 8pt,
  frame           = single,
  framerule       = 0.4pt,
  rulecolor       = \color{secondary!40},
  breaklines      = true,
  captionpos      = b,
  tabsize         = 2,
}

% --- Matematica ---
\usepackage{amsmath, amssymb}

% --- Enumerazioni personalizzate ---
\usepackage{enumitem}
\setlist[itemize]{leftmargin=1.5em, topsep=4pt, itemsep=2pt}
\setlist[enumerate]{leftmargin=1.5em, topsep=4pt, itemsep=2pt}

% --- Caption ---
\usepackage{caption}
\captionsetup{
  font        = small,
  labelfont   = {bf, color=primary},
  labelsep    = period,
  skip        = 6pt,
}

% --- Spaziatura ---
\usepackage{setspace}
\setstretch{1.15}
\setlength{\parskip}{6pt}
\setlength{\parindent}{0pt}

% --- Titolo sezione personalizzato ---
\usepackage{titlesec}
\titleformat{\section}
  {\large\bfseries\color{primary}}
  {\thesection.}{0.6em}{}[\titlerule]
\titleformat{\subsection}
  {\normalsize\bfseries\color{primary!80}}
  {\thesubsection.}{0.5em}{}
\titleformat{\subsubsection}
  {\normalsize\itshape\color{secondary}}
  {\thesubsubsection.}{0.5em}{}

\titlespacing{\section}{0pt}{14pt}{6pt}
\titlespacing{\subsection}{0pt}{10pt}{4pt}

% --- Riquadro informativo ---
\usepackage{tcolorbox}
\tcbuselibrary{skins, breakable}
\newtcolorbox{infobox}[1][]{
  enhanced,
  colback  = primary!5,
  colframe = primary!50,
  fonttitle= \bfseries,
  title    = {#1},
  breakable,
}

% --- Utility ---
\usepackage{lipsum}   % testo segnaposto — rimuovere nella versione finale
\usepackage{xspace}
\newcommand{\mvc}{\textsc{mvc}\xspace}
\newcommand{\teamid}[1]{\textbf{#1}}
